\documentclass[dvisvgm,tikz,border=3pt]{standalone}
\usepackage{amsmath,amssymb}
\usepackage{pgfplots}
\usepackage{CJKutf8}
\pgfplotsset{compat=1.18}
\usetikzlibrary{arrows.meta,calc}

\definecolor{themebg}{HTML}{2e362c}
\definecolor{themefg}{HTML}{edf4e8}
\definecolor{themegray}{HTML}{9ea897}
\definecolor{themecurve}{HTML}{7eb6ff}
\definecolor{themealert}{HTML}{ff7b7b}
\definecolor{themeamber}{HTML}{f5b942}

\begin{document}
\begin{CJK*}{UTF8}{gbsn}
\pagecolor{themebg}
\begin{tikzpicture}[color=themefg, text=themefg]

% ==================== Plot 1: Two Axes (T = 2d) ====================
\begin{axis}[
    name=ax1,
    width=7.2cm,
    height=6.2cm,
    axis lines=middle,
    axis line style={color=themefg, -{Stealth}},
    tick style={color=themefg},
    tick label style={color=themefg, font=\small},
    every axis label/.style={color=themefg},
    xmin=-1.2, xmax=4.4,
    ymin=-1.4, ymax=2.2,
    xtick={-1, 0, 1, 2, 3, 4},
    xticklabels={$-d$, $0$, $d$, $2d$, $3d$, $4d$},
    ytick={0},
    clip=false,
    xlabel={$x$},
    ylabel={$y$},
    title style={at={(0.5,1.12)},anchor=south,color=themefg,font=\bfseries\small},
    title={双轴对称（$x=0,\;x=d$）},
]
  % Axes x=0 and x=d
  \draw[dashed, themegray, line width=1.1pt] (axis cs:0, -1.1) -- (axis cs:0, 1.8);
  \draw[dashed, themegray, line width=1.1pt] (axis cs:1, -1.1) -- (axis cs:1, 1.8);

  % 双轴示意必须主动排除“d 也像周期”的错觉。
  % 取严格递增且端点不等的非对称母段 g_1(u)=0.25+1.25u-0.5u^2+0.1u^3，u∈[0,1]。
  % f(0)=0.25, f(d)=1.10，因此 d 明确不是该示意函数的周期；其余区间只由两条轴真实镜像生成。
  \addplot[themegray, line width=1.4pt, dashed, domain=-1:0, samples=100] {0.25-1.25*x-0.5*x^2-0.1*x^3};
  \addplot[themecurve, line width=2pt, domain=0:1, samples=100] {0.25+1.25*x-0.5*x^2+0.1*x^3};
  \addplot[themecurve, line width=2pt, domain=1:2, samples=100] {0.25+1.25*(2-x)-0.5*(2-x)^2+0.1*(2-x)^3};
  \addplot[themegray, line width=1.4pt, dashed, domain=2:3, samples=100] {0.25+1.25*(x-2)-0.5*(x-2)^2+0.1*(x-2)^3};
  \addplot[themegray, line width=1.4pt, dashed, domain=3:4, samples=100] {0.25+1.25*(4-x)-0.5*(4-x)^2+0.1*(4-x)^3};
  \addplot[only marks, mark=*, mark size=2.4pt, fill=themecurve, draw=themefg] coordinates {(0,0.25) (1,1.10)};

  % Period Bracket
  \draw[<->, line width=1.2pt, color=themecurve] (axis cs:0, 1.5) -- node[above, font=\small\bfseries] {一个周期 $2d$} (axis cs:2, 1.5);
\end{axis}

\node[font=\small\bfseries, color=themecurve, anchor=north] at ($(ax1.south)-(0, 0.45cm)$) {$f(0)\ne f(d)\Rightarrow d\text{ 不是周期};\qquad 2d\text{ 是一个周期}$};

% ==================== Plot 2: Two Centers (T = 2d) ====================
\begin{axis}[
    at={($(ax1.east)+(2.0cm,0)$)},
    anchor=west,
    name=ax2,
    width=7.2cm,
    height=6.2cm,
    axis lines=middle,
    axis line style={color=themefg, -{Stealth}},
    tick style={color=themefg},
    tick label style={color=themefg, font=\small},
    every axis label/.style={color=themefg},
    xmin=-1.2, xmax=4.4,
    ymin=-1.6, ymax=2.0,
    xtick={-1, 0, 1, 2, 3, 4},
    xticklabels={$-d$, $0$, $d$, $2d$, $3d$, $4d$},
    ytick={0},
    clip=false,
    xlabel={$x$},
    ylabel={$y$},
    title style={at={(0.5,1.12)},anchor=south,color=themefg,font=\bfseries\small},
    title={双中心对称（$(0,0),\;(d,0)$）},
]
  % Center points
  \fill[themefg] (axis cs:0, 0) circle (2.8pt);
  \fill[themefg] (axis cs:1, 0) circle (2.8pt);

  % 两中心示意采用更强非对称母段 g_2(u)=u(1-u)(1+2u)。
  % 两侧严格按中心对称 f(c-t)=-f(c+t) 生成；平移 d 后正负与左右形态均不重合。
  \addplot[themegray, line width=1.4pt, dashed, domain=-1:0, samples=100] {-((-x)*(1+x)*(1-2*x))};
  \addplot[themecurve, line width=2pt, domain=0:1, samples=100] {x*(1-x)*(1+2*x)};
  \addplot[themecurve, line width=2pt, domain=1:2, samples=100] {-((2-x)*(x-1)*(1+2*(2-x)))};
  \addplot[themegray, line width=1.4pt, dashed, domain=2:3, samples=100] {(x-2)*(3-x)*(1+2*(x-2))};
  \addplot[themegray, line width=1.4pt, dashed, domain=3:4, samples=100] {-((4-x)*(x-3)*(1+2*(4-x)))};

  % Period Bracket
  \draw[<->, line width=1.2pt, color=themecurve] (axis cs:0, 1.4) -- node[above, font=\small\bfseries] {一个周期 $2d$} (axis cs:2, 1.4);
\end{axis}

\node[font=\small\bfseries, color=themecurve, anchor=north] at ($(ax2.south)-(0, 0.45cm)$) {$f(\frac d4)\ne f(\frac{5d}{4})\Rightarrow d\text{ 不是周期};\qquad 2d\text{ 是一个周期}$};

% ==================== Plot 3: One Axis + One Center (T = 4d) ====================
\begin{axis}[
    at={($(ax1.south west)!0.5!(ax2.south east)-(0,2.7cm)$)},
    anchor=north,
    name=ax3,
    width=7.2cm,
    height=6.2cm,
    axis lines=middle,
    axis line style={color=themefg, -{Stealth}},
    tick style={color=themefg},
    tick label style={color=themefg, font=\small},
    every axis label/.style={color=themefg},
    xmin=-1.2, xmax=4.4,
    ymin=-1.6, ymax=2.0,
    xtick={-1, 0, 1, 2, 3, 4},
    xticklabels={$-d$, $0$, $d$, $2d$, $3d$, $4d$},
    ytick={0},
    clip=false,
    xlabel={$x$},
    ylabel={$y$},
    title style={at={(0.5,1.12)},anchor=south,color=themefg,font=\bfseries\small},
    title={一轴一中心（轴 $x=0$，中心 $(d,0)$）},
]
  % Axis x=0 and Center (1,0)
  \draw[dashed, themegray, line width=1.1pt] (axis cs:0, -1.2) -- (axis cs:0, 1.8);
  \fill[themefg] (axis cs:1, 0) circle (2.8pt);

  % 轴 x=0 给出 f(-x)=f(x)；中心 (1,0) 给出 f(1-t)=-f(1+t)。
  % 取 g_3(u)=0.6+0.7u-1.3u^2，使 f(0)=0.6 而 f(d)=0，直接排除 d 周期；
  % 两结构复合得到 f(x+2d)=-f(x)，且 f(0)≠0，所以 2d 也不是该示意函数的周期。
  \addplot[themegray, line width=1.4pt, dashed, domain=-1:0, samples=100] {0.6-0.7*x-1.3*x^2};
  \addplot[themealert, line width=2pt, domain=0:1, samples=100] {0.6+0.7*x-1.3*x^2};
  \addplot[themealert, line width=2pt, domain=1:2, samples=100] {-(0.6+0.7*(2-x)-1.3*(2-x)^2)};
  \addplot[themealert, line width=2pt, domain=2:3, samples=100] {-(0.6+0.7*(x-2)-1.3*(x-2)^2)};
  \addplot[themealert, line width=2pt, domain=3:4, samples=100] {0.6+0.7*(4-x)-1.3*(4-x)^2};

  % Half period & Full period
  \draw[<->, line width=1.1pt, color=themegray] (axis cs:0, -1.3) -- node[below, font=\scriptsize] {$2d\text{ 后反号}$} (axis cs:2, -1.3);
  \draw[<->, line width=1.3pt, color=themealert] (axis cs:0, 1.4) -- node[above, font=\small\bfseries] {一个周期 $4d$} (axis cs:4, 1.4);
\end{axis}

\node[font=\small\bfseries, color=themealert, anchor=north] at ($(ax3.south)-(0, 0.45cm)$) {$d,2d\text{ 均不是周期};\qquad 4d\text{ 是一个周期}$};

\end{tikzpicture}
\end{CJK*}
\end{document}
